\documentclass[11pt]{article}
\usepackage[margin=1in]{geometry}
\usepackage{amsmath,amssymb,amsthm}
\newtheorem{theorem}{Theorem}
\title{An exact atom exchange and its optimal barrier for four fair bits}
\author{Joshua Mark}
\date{20 November 2026; research draft, priority unestablished}
\begin{document}\maketitle
For four unbiased pairwise-independent sign bits, the free Walsh
coefficients are $a\in\mathbb R^5$, comprising the four triple
characters and the quadruple character. The feasible polytope is
\[
 P=\{a:1+v\cdot a\ge0\text{ for all }v\in\{\pm1\}^5,
 \ \prod_i v_i=1\}.
\]
The map from bit strings to these sixteen $v$ is a bijection.
The previously established vertex description is
$\{\pm e_i\}_{i=1}^5\cup\{s/3:\prod_i s_i=1\}$.
Its collision probability is $C(a)=(1+\|a\|^2)/16$.
The ten first vertices give uniform eight-atom laws with collision
$1/8$; the sixteen latter vertices have one atom of weight $1/6$,
ten of weight $1/12$, and collision $7/72$.

\begin{theorem}[A finite exact escape rule]
At any weighted vertex $a=s/3$, choose any $i\in\{1,\ldots,5\}$
and replace the law by the uniform vertex $b=s_i e_i$.
This is an edge move of the feasible moment polytope, preserves all
normalization, positivity and pairwise constraints, and changes
twelve atom weights. It increases collision by exactly $1/36$.
Along its feasible interpolation $a(t)=(1-t)a+tb$,
\[
 C(a(t))=\frac7{72}-\frac{t}{36}+\frac{t^2}{18}.
\]
It initially decreases collision, returns to the initial value at
$t=1/2$, and reaches its minimum $3/32$ at $t=1/4$.
\end{theorem}
\begin{proof}
The weighted vertex has five active inequalities, with normals
$s$ multiplied coordinatewise by each sign vector having four minus
signs. At $b$, all except the normal whose plus position is $i$
remain active. These four normals are independent, so their common
face has dimension at most one. The segment between the two distinct
vertices lies in that face, hence is an edge. Its midpoint is a
strict convex combination of two laws; no nonnegativity is lost.
The weighted support has eleven atoms, the uniform support eight,
and their intersection seven, leaving a twelve-atom union.
All twelve weights change: the nonzero differences have magnitudes
$1/24$, $1/12$ or $1/8$. Expanding
$\|a(t)\|^2=5/9-4t/9+8t^2/9$ proves the formula.
The edge is a circuit in the normalization/pairwise moment kernel:
its support-restricted kernel has dimension one, and every entry of
its direction on those twelve atoms is nonzero.
\end{proof}

\begin{theorem}[The unavoidable barrier is exactly $1/288$]
Every continuous feasible path from a weighted vertex to any uniform
vertex passes through a law of collision at most $3/32$.
The preceding edge path attains this upper bound on its minimum, so
it is optimal in the maximin sense. In particular no continuously
collision-nondecreasing path escapes a weighted vertex to an optimum.
\end{theorem}
\begin{proof}
We first show that every $a\in P$ with $|a_i|=1/2$ satisfies
$\|a\|^2\le1/2$. Permute coordinates to take $i=1$.
If $a_1=1/2$, write $y=(a_2,\ldots,a_5)$. The constraints with
$v_1=-1$ use the eight four-dimensional odd-parity sign vectors.
They consist of $\pm h_j$, where the four rows $h_j$ of $J_4-2I_4$
are mutually orthogonal with squared norm four. Thus
$|h_j\cdot y|\le1/2$, and
\[
 \|y\|^2=\frac14\sum_{j=1}^4(h_j\cdot y)^2\le1/4.
\]
If $a_1=-1/2$, use $v_1=+1$ instead; the even-parity four-sign
vectors likewise consist of the signed rows of a Hadamard matrix
of order four. The identical bound follows. More generally the
argument shows $\|a\|^2\le a_i^2+(1-|a_i|)^2$.

At a weighted vertex all five coordinate magnitudes are $1/3$;
at a uniform vertex one is one. The intermediate value theorem
forces some coordinate to cross magnitude $1/2$ along every
continuous connecting path. At that point collision is at most
$(1+1/2)/16=3/32$. The displayed edge path has minimum exactly that
value. Since $7/72-3/32=1/288$, the claimed barrier follows.
\end{proof}

This is an actual nonlocal twelve-atom correction rule on the fully
solved four-bit testbed. It demonstrates that strict local collision
maxima can require a quantified loss before an exact feasible path
reaches a better distribution. It does not assert that every higher
dimensional moment polytope has such a bounded circuit or favorable
endpoint. The four-dimensional Hadamard frame used to certify the
barrier is supplied explicitly; no unknown-order construction is
assumed or obtained.
\end{document}
